\section{Introduction}
\label{sec:intro}

Language models now summarize novels, write screenplays, and answer questions about complex narratives---all tasks that require keeping track of who is where and who has what.
But how many characters can a model actually track at once?
Prior work on entity tracking has either fixed the entity count at a single value~\citep{kim2023entity} or used naturalistic texts where character count and document length are hopelessly confounded~\citep{hamilton2025tldm}.
No study has systematically varied the number of tracked entities while holding other factors constant.

We fill this gap.
We construct synthetic narratives containing 2 to 64 named characters, each with a location and a possession that change over the course of the story.
After presenting the full narrative, we ask the model to report every character's final state via structured JSON output.
By comparing model responses against programmatic ground truth, we measure accuracy as a function of character count and classify every error as a \emph{confusion} (the model's answer belongs to a different character), a \emph{forgetting} error (the answer reverts to the character's initial state), or \emph{other}.

Our main findings are:
\begin{itemize}[leftmargin=*,itemsep=0pt,topsep=0pt]
\item We show that frontier 2025 models (\gptfull and \gptmini) achieve \textbf{perfect accuracy up to 16~characters} and maintain $>$93\% accuracy at 64~characters---substantially higher than prior benchmarks suggested (\eg 65\% for Code Llama 70B on 7~entities;~\citealt{kim2024code}).
\item We find that errors are overwhelmingly \textbf{confusions} (93.9\% for \gptmini), not forgetting, implying that character representations overlap rather than occupying independent memory slots.
\item We observe \textbf{gradual, not cliff-like} degradation, with no sharp transition point---suggesting a continuous capacity mechanism rather than a discrete slot limit.
\item We discover that \gptmini \textbf{outperforms} \gptfull, and identify a specific failure mode in \gptfull where compound drop-then-pickup sentences cause possession-tracking errors.
\end{itemize}

These results bear on both practice and theory.
For practitioners, they establish that frontier models are reliable for narratives with up to $\sim$16 characters without augmentation; beyond 32, retrieval-augmented approaches may help.
For researchers studying transformer internals, the confusion-dominated error pattern provides behavioral evidence consistent with the associative tracking algorithms identified by \citet{li2025state} through mechanistic interpretability.

The rest of this paper is organized as follows.
\Secref{sec:related} surveys prior work on entity tracking and narrative understanding.
\Secref{sec:method} describes our synthetic story generation, experimental protocol, and error taxonomy.
\Secref{sec:results} presents the main results, and \secref{sec:discussion} interprets them.
\Secref{sec:conclusion} concludes.
